\documentclass[11pt]{article}

\usepackage[T1]{fontenc}
\usepackage{amsmath,amssymb,amsthm}
\usepackage{hyperref}

\title{AIC, TIC, and RMSEA as Bias Corrections}
\author{}
\date{\today}

\newcommand{\E}{\mathbb{E}}
\newcommand{\tr}{\operatorname{tr}}
\newcommand{\argmax}{\operatorname*{arg\,max}}
\newcommand{\argmin}{\operatorname*{arg\,min}}
\newcommand{\KL}{\operatorname{KL}}
\newcommand{\dto}{\xrightarrow{d}}
\newcommand{\pto}{\xrightarrow{p}}

\theoremstyle{plain}
\newtheorem{proposition}{Proposition}
\theoremstyle{remark}
\newtheorem*{remark}{Remark}

\begin{document}
\maketitle

\section*{Purpose}

This note records the common shape behind AIC/TIC and RMSEA. Both start from a
sample criterion that is too favorable because it is evaluated after fitting.
Both subtract a first-order bias term. The bias terms are different because the
targets are different:
\[
  \text{AIC/TIC: predictive risk of the fitted distribution,}
\]
\[
  \text{RMSEA: population model misfit, or noncentrality per observation per df.}
\]
AIC and TIC do not require the exact-fit or local-alternative story that often
motivates RMSEA. They use local Taylor expansions around the fitted parameter
and its pseudo-true target, but the statistical target is fixed-misspecification
predictive risk, not a Pitman drift from a correct model.

\section{AIC: correcting in-sample optimism}

Let $Y_1,\ldots,Y_n\sim g$ and let $f_\theta$ be a parametric working model.
Write the log likelihood
\[
  \ell_n(\theta)=\sum_{i=1}^n \log f_\theta(Y_i),
  \qquad
  \hat\theta=\argmax_\theta \ell_n(\theta).
\]
The population target of maximum likelihood, even when the model is wrong, is
the KL projection
\[
  \theta_0=\argmax_\theta M(\theta),
  \qquad
  M(\theta)=\E_g\{\log f_\theta(Y)\}.
\]
Equivalently, $\theta_0$ minimizes $\KL(g,f_\theta)$ over the model class.

The quantity we would like to estimate is the expected log score of the fitted
model on a fresh observation,
\[
  R_n = \E_{\hat\theta}\E_{Y^0}\{\log f_{\hat\theta}(Y^0)\},
\]
where $Y^0\sim g$ is independent of the training sample and the outer
expectation averages over the training sample. The empirical maximized log
likelihood $n^{-1}\ell_n(\hat\theta)$ estimates the same log score on the
training observations, not on fresh observations. It is optimistic because
$\hat\theta$ was chosen to fit those same observations.

\begin{proposition}[Akaike's optimism correction]
If the model is correctly specified, regular, and has $k$ free parameters, then
\[
  \E\{\ell_n(\hat\theta)\}
  =
  n\,\E\{\log f_{\hat\theta}(Y^0)\} + k + o(1).
\]
Therefore
\[
  \ell_n(\hat\theta)-k
\]
is an asymptotically unbiased estimator, up to constants independent of the
candidate model, of the expected out-of-sample log likelihood. Multiplying by
$-2$ gives
\[
  \mathrm{AIC}=-2\ell_n(\hat\theta)+2k .
\]
\end{proposition}

\begin{proof}[Derivation sketch]
Let
\[
  s_i(\theta)=\frac{\partial\log f_\theta(Y_i)}{\partial\theta},
  \qquad
  I=-\E\left\{
    \frac{\partial^2\log f_{\theta_0}(Y)}
         {\partial\theta\partial\theta^\top}
  \right\}.
\]
Under correct specification,
$\sqrt n(\hat\theta-\theta_0)\dto N(0,I^{-1})$ and
$\E\{s_i(\theta_0)s_i(\theta_0)^\top\}=I$.

For a fresh observation,
\[
  M(\hat\theta)
  =
  M(\theta_0)
  -\frac12(\hat\theta-\theta_0)^\top I(\hat\theta-\theta_0)
  +o_p(n^{-1}),
\]
so
\[
  n\,\E M(\hat\theta)=nM(\theta_0)-\frac12 k+o(1).
\]
For the training log likelihood, the score term does not average away because
$\hat\theta$ is a function of the same data. A quadratic expansion of
$\ell_n(\hat\theta)$ around $\theta_0$ gives
\[
  \E\ell_n(\hat\theta)=nM(\theta_0)+\frac12 k+o(1).
\]
The gap is $k+o(1)$ on the log-likelihood scale, or $2k+o(1)$ on the deviance
scale. This is the AIC penalty.
\end{proof}

\begin{remark}
AIC is often described as making the maximized log likelihood approximately
unbiased for expected predictive log likelihood. More precisely, it removes
the first-order \emph{optimism} of reusing the training data for both fitting
and evaluation. It is not a small-sample unbiasedness claim, and it is not a
test of exact fit.
\end{remark}

\section{TIC: the same correction under misspecification}

When the model is misspecified, the Hessian and score variance no longer agree.
Define
\[
  I = -\E_g\left\{\frac{\partial^2\log f_{\theta_0}(Y)}
    {\partial\theta\partial\theta^\top}\right\},
  \qquad
  J = \E_g\{s(Y;\theta_0)s(Y;\theta_0)^\top\}.
\]
Then
\[
  \sqrt n(\hat\theta-\theta_0)\dto N(0,I^{-1}JI^{-1}).
\]
The same local expansion gives the optimism
\[
  \tr(I^{-1}J)
\]
on the log-likelihood scale. Takeuchi's information criterion is therefore
\[
  \mathrm{TIC}
  =
  -2\ell_n(\hat\theta)
  +2\,\tr(\hat I^{-1}\hat J).
\]
Correct specification is the special case $I=J$, giving
$\tr(I^{-1}J)=k$ and recovering AIC.

\begin{remark}
This is the sense in which AIC/TIC use local arguments indirectly. The target is
fixed-population predictive risk, but the bias calculation is a second-order
expansion around the KL projection. No local drift toward a true finite
dimensional model is needed.
\end{remark}

\section{RMSEA: correcting chi-square noise}

Now let $\hat\eta$ be a consistent estimator of a saturated population object
$\eta$, such as $(\mu,\Sigma)$ or a limited-information ordinal moment vector.
Let $\mathcal C\subset E$ be the model-implied class and let
\[
  F_{\mathcal C}(\eta)
  =
  \inf_{c\in\mathcal C} D(\eta,c)
\]
be a nonnegative discrepancy, with $F_{\mathcal C}(\eta)=0$ iff
$\eta\in\mathcal C$. Continuity gives
\[
  F_{\mathcal C}(\hat\eta)\pto F_{\mathcal C}(\eta)
\]
even when $\hat\eta$ has finite-sample bias. The population RMSEA target is
\[
  \varepsilon_{\mathcal C}(\eta)
  =
  \sqrt{\frac{F_{\mathcal C}(\eta)}{df}}.
\]

The usual SEM test statistic has the form
\[
  T_n = n F_{\mathcal C}(\hat\eta)
\]
after fitting the projection onto $\mathcal C$; in magmaan's convention
$T_n=2n\,\texttt{fmin}$. Under exact fit and the usual regularity conditions,
\[
  T_n \dto \chi^2_{df},
\]
so pure sampling noise contributes about $df$ to the statistic even when the
population discrepancy is zero. Under fixed misspecification,
\[
  T_n = nF_{\mathcal C}(\eta)+O_p(\sqrt n),
\]
so the exact-fit test diverges, but the rescaled discrepancy remains meaningful.
The RMSEA plug-in correction is
\[
  \widehat\varepsilon
  =
  \sqrt{
    \frac{\max(T_n-df,0)}{n\,df}
  }.
\]

\begin{remark}
The subtraction of $df$ in RMSEA is analogous in spirit to the AIC subtraction
of $k$: both remove a first-order fitting/noise contribution. But the targets
are not the same. AIC/TIC correct an in-sample log score to estimate
out-of-sample predictive risk. RMSEA corrects a chi-square-like lack-of-fit
statistic to estimate excess noncentrality, hence population discrepancy per
degree of freedom.
\end{remark}

\section{Where local alternatives enter}

The population definition
\[
  \varepsilon_{\mathcal C}(\eta)=\sqrt{F_{\mathcal C}(\eta)/df}
\]
does not require local alternatives. It is a fixed-population estimand. Local
alternatives enter when one wants the chi-square reference distribution,
confidence intervals, close-fit tests, or power calculations to remain
nondegenerate. A Pitman sequence writes
\[
  \eta_n=\eta_0+h/\sqrt n,\qquad \eta_0\in\mathcal C,
\]
so
\[
  nF_{\mathcal C}(\eta_n)\to\lambda
  \quad\text{and}\quad
  T_n\dto\chi^2_{df}(\lambda).
\]
Then RMSEA is the square root of $\lambda/(n\,df)$ on the sample scale, or of
$F_{\mathcal C}(\eta_n)/df$ on the population-discrepancy scale. Without the
local sequence, fixed positive misfit makes $T_n$ diverge, so there is no
finite noncentral chi-square limit to invert; there is still a perfectly good
population discrepancy.

\section{Comparison}

\begin{center}
\begin{tabular}{p{0.14\linewidth}p{0.52\linewidth}p{0.24\linewidth}}
criterion & target & bias/noise correction\\
\hline
AIC & predictive log likelihood, correct model & $k$ on log-likelihood scale\\
TIC & predictive log likelihood, misspecified model & $\tr(I^{-1}J)$\\
RMSEA & population discrepancy per df & $df$ chi-square noise
\end{tabular}
\end{center}

Thus the closest one-line memory aid is:
\begin{align*}
  &\text{AIC/TIC estimate prediction risk after fitting;}\\
  &\text{RMSEA estimates lack of fit after removing expected test noise.}
\end{align*}

\end{document}
